% =====================================================================
%  Chapter 12 — Step 6
% =====================================================================

\chapter{Step Six: The Rallis Seesaw and the Bessel Period}
\label{ch:seesaw}

The argument needs to average over the Hecke characters, and it cannot do that
directly, because the object it has --- a period integral attached to a
character --- depends on the character in a way that makes the average
useless. The Rallis seesaw is what removes the dependence.

The remark's item 12:

\begin{quote}
\small
``Seesaw: the period of the iterated lift is a Bessel period independent of the
Hecke characters, Rallis.''
\end{quote}

One sentence, and it carries a classical piece of machinery.

\section{What the seesaw is}The seesaw principle, in the form Rallis uses it, says that a period integral
of a Siegel--Weil theta series on a unitary group is equal, up to explicit
factors, to a period integral of the \emph{restriction of scalars} of that
group to a definite unitary group, matched against an Eisenstein series.

Concretely: if $G$ is a unitary group over a CM field and $G^0$ is the
totally definite unitary group obtained by restricting scalars to the maximal
real subfield, then for suitable Eisenstein series $E_1$ on $G$ and $E_2$ on
$G^0$,
\begin{equation}
\int_{G^0\backslash G} \theta_A(g)\,E_1(g)\,dg
\;=\;
\frac{\hat g \cdot L_A(1/2,\varepsilon)}{L(1,\varepsilon)}\int_{G\backslash G} \theta_A(g)\,E_2(g)\,dg,
\label{eq:seesaw}
\end{equation}
where $\hat g$ is the number of components of $G^0\backslash G$ and $L_A$ is the
$L$-function of the form $A$.

The point of the formula is not the equality. The point is the \emph{structure}
of the right-hand side: the integral on $G^0\backslash G$ is a period over a
compact quotient, computed with an Eisenstein series --- and Eisenstein series
are attached to \emph{Hecke characters}. The left-hand side has the character
sitting in the integral, where it is hard to average over; the right-hand side
has it sitting in $L_A(1/2,\varepsilon)$, where it is an $L$-value and can be
averaged.

\begin{quote}
The seesaw is what lets the argument average. Before it, the character is in
the way of the integral; after it, the character has been traded for an
$L$-value, and $L$-values average.
\end{quote}

\section{Why the period becomes a Bessel period}Apply the seesaw to the theta lift. A theta lift to a group of higher rank
has a period that is, by the structure of the theta series, a product of
ordinary Bessel functions: the Weil representation restricted to a torus gives
a one-dimensional character whose $L$-function is a product of Bessel
$L$-functions, and the theta series is essentially a sum over that character.

Concretely, for the lift of $\pi$ from $\mathrm{U}(E,n)$ to $\mathrm{U}(E,d)$,
the period is
\begin{equation}
B_A(\phi, F) \;=\; \sum_{\gamma} \phi(\gamma)\,\theta_A(\gamma)\,F(\gamma)
\label{eq:bessel}
\end{equation}
for a suitable $\phi$ and a cusp form $F$, and the point is that this
quantity factors into Bessel functions and does not depend on the Hecke
character $\mu$ except through a fixed character that has already been
extracted.

The remark's phrase ``independent of the Hecke characters'' should be read
carefully. It means: the period that the argument needs to average over the
Hecke characters does not itself carry the Hecke character. The character has
been absorbed into the summand of the theta series, via the notation
$\pi \mapsto \pi \tboxplus \mu_j[t]$ in the remark's item 8: the $\mu_j$ appear
in the \emph{index} of the lift, not in the period.

That distinction is exactly what makes the rest of the argument possible. Once
the period is character-independent, the summation over characters becomes a
summation over a fixed function, which the orbit-sum identity of
Chapter~\ref{ch:step4} can decompose.

\section{The formalisation of the descent tower}The development states all of this. In
\code{Formalism.SharpRestriction.Divisibility.FunctionModelRallisFromInputs.lean}
the relevant definitions are:

\begin{seamnote}{The Bessel period and the lift operator}
{\footnotesize
\begin{leanquote}
def bB {D : D7} (X : ThetaX D) (t j : Nat)
    (pos : X.Pos (X.Xs t j) (X.Ys t j) (X.iXYs t j))
(s : X.FJd (X.Xs t j) (X.Ys t j) (X.iXYs t j)) (phi : X.AF
      (X.Xs t j))
    (F : X.cusp (X.Ys t j)) : Complex :=
  X.Bj t j pos s (X.conjAF _ phi) F

def vth {D : D7} (X : ThetaX D) (t : Nat) (mu : Nat ->
  X.Chr) (j : Nat)
(Phi : X.Sch (X.Ys t (j + 1)) (X.Ys t j)) : X.cusp (X.Ys t
      (j + 1)) ->L[Complex] X.AF (X.Ys t j) :=
X.th _ _ (car X mu (j + 1) (j + 1)) (car X mu (j + 1) j) Phi
\end{leanquote}
}
\end{seamnote}

Read them together and they are the mathematics of the step.

\begin{description}[style=nextline,leftmargin=0pt,font=\sffamily\bfseries\small,
  itemsep=6pt]
\item[\code{bB}] is the Bessel period: a pairing of the conjugate of an
automorphic form $\varphi$ with a cusp form $F$, evaluated through the
Bessel-type operator \code{Bj} at the place $\mathrm{pos}$. It takes the
\emph{cuspidal form} $F$ and the \emph{automorphic form} $\varphi$ --- the two
objects the seesaw brings together.
\item[\code{vth}] is the theta lift operator itself: from the space of
cuspidal forms at stage $j+1$ to the space of automorphic forms at stage $j$,
via \code{X.th} with the character $\mathrm{car}\ X\ \mu\ (j+1)\ (j+1)$, which
is $\mu_{j+1}$ or $\mu_{j+1}\eta^{-1}$ according to the parity of the index.
\end{description}

And the conditions that must hold at each stage are:

\begin{seamnote}{\code{PerNZ} --- the period must not vanish}
{\footnotesize
\begin{leanquote}
def PerNZ {D : D7} (X : ThetaX D) (t j : Nat) (phi : X.AF
  (X.Xs t j)) (F : X.cusp (X.Ys t j)) : Prop :=
  (j < D.R.m and exists pos s, bB X t j pos s phi F != 0) or
exists h : j = D.R.m, X.Btop (X.Xs t j) (X.Ys t j) (h >
      X.iYXm t) (X.conjAF _ phi) F != 0
\end{leanquote}
}
\end{seamnote}

That is exactly the right condition and it is worth pausing on. The Bessel
period \emph{could} vanish. If it does, the seesaw gives zero on one side, and
the argument has nothing to work with. So the theorem must establish that it
does not vanish --- either at an interior stage $j < m$, or at the top stage
$j = m$ via the top-level operator \code{Btop}. The condition is a
disjunction, and the disjunction is doing real work: some stages are handled
by the ordinary Bessel operator and the last stage by the top-level one.

\begin{seamnote}{\code{IsAutRep} --- the lifted form is genuinely automorphic}
{\footnotesize
\begin{leanquote}
def IsAutRep {D : D7} (X : ThetaX D) (V : TH.HSp D) (pi :
  Submodule Complex (X.AF V)) : Prop :=
  X.IsStable V pi and pi != bottom and
forall pi' : Submodule Complex (X.AF V), pi' <= pi ->
      X.IsStable V pi' -> pi' = bottom or pi' = pi
\end{leanquote}
}
\end{seamnote}

The second clause of this is the irreducibility condition, and it is what
distinguishes a cuspidal lift from an Eisenstein one. A theta lift of a cusp
form is usually irreducible; a theta lift of an Eisenstein series is not, and
its period is a product of $L$-values rather than a new $L$-function. The
method needs the cuspidal case throughout, because the whole argument is
about the $L$-function of the zeta function and the Weil expression built
out of its zeros.

So the descent-tower conditions are, in order of what they check:

\begin{center}
\begin{tabular}{@{}l>{\RaggedRight\arraybackslash\hyphenpenalty=10000\exhyphenpenalty=10000}p{0.68\textwidth}@{}}
\toprule
\textbf{Condition} & \textbf{What it guarantees} \\
\midrule
\code{IsDv} & The type condition is preserved at each stage. \\
\code{IsAutRep} & The lifted representation is irreducible, hence cuspidal. \\
\code{PerNZ} & The Bessel period does not vanish, so the seesaw carries
information. \\
\code{IsMu} & The characters $\mu_j$ are chosen so that the weight inequality
at each stage holds. \\
\code{Hyp3} & The global hypotheses: interlacing, parity, centrality, and the
cubic gap condition. \\
\bottomrule
\end{tabular}
\end{center}

\section{The parity, and how the character is extracted}One detail in the formalisation is worth a reader's attention, because it is
the mechanism by which the character is made to disappear. Recall:

\begin{seamnote}{\code{car}}
{\footnotesize
\begin{leanquote}
def car {D : D7} (X : ThetaX D) (mu : Nat -> X.Chr) (s i :
  Nat) : X.Chr :=
  if Even i then mu s * X.etaT⁻¹ else mu s
\end{leanquote}
}
\end{seamnote}

The character used at stage $i$ is $\mu_s$ for even $i$ and $\mu_s\eta^{-1}$ for
odd $i$, where $\eta$ is the fixed character attached to the tower. So the
tower alternates between $\mu$ and $\mu\eta^{-1}$.

\begin{quote}
\textbf{Why this matters.} The alternation is what makes the parity condition
\code{TParity} satisfiable at all: the weight of a theta series of rank $n$ is
congruent to $n \bmod 2$, so the weights available at successive stages
alternate in parity, and the characters attached to them must alternate with
them. The condition \code{IsMu} says precisely that
$X.\mathrm{infTy}(\mu_{j+1}) = -P.e(\tau, j)$ --- the infinity type of the
character at stage $j+1$ is the negative of the parameter $e_j$, which is the
weight-matching condition of \eqref{eq:weightineq}.
\end{quote}

So the tower is: a fixed character $\eta$ times the user's character $\mu$,
alternating, with the weight at each stage matching the form at that stage.

\section{What step 6 buys}The seesaw converts a character-dependent period into an $L$-value, and the
Bessel structure of the theta series converts the summand into something that
does not carry the character at all. Between them, the character is available
to be averaged over, which is what the rest of the argument needs. The cost
is the condition \code{PerNZ}: the period must not vanish, and that is a
non-trivial thing to prove.

\begin{keytakeaways}
\item The seesaw trades a character-dependent period integral for a period
against an Eisenstein series on the totally definite group, and the character
ends up in an $L$-value.
\item The theta lift's period is a product of Bessel functions, so it does
not carry the Hecke character; the character is absorbed into the index of
the lift, as $\pi \mapsto \pi \tboxplus \mu_j[t]$.
\item The price is non-vanishing: \code{PerNZ} requires the Bessel period to
be non-zero, at an interior stage or at the top stage, and this is a genuine
obligation.
\item The lifted form must stay irreducible (\code{IsAutRep}), or the period
degenerates to an $L$-value and the argument has nothing to work with.
\end{keytakeaways}
